\documentclass[book.tex]{subfiles}

\begin{document}

\chapter{Introduction to Type Theory}

\label{chap:type_theory}
\noindent\rule{11cm}{0.4pt} \\
\textbf{Prerequisites:}  \autoref{chap:math_basics},   \\
\textbf{Difficulty Level:} ** \\
\noindent\rule{11cm}{0.4pt}
\vspace{5mm}


Type theory is an increasingly important part of modern mathematics and computer science that provides technology to describe the behavior of computer programs.. In the rest of this book, we will use types to help us more accurate characterize physical systems.

\section{Types, Terms, and Values}

The parts of a program are called terms, types and values. Terms are program statements and expressions such as 

\begin{lstlisting}
3+4, x, 234
\end{lstlisting}

A term is usually denoted with the symbol $e$. Each term is associated with a type that provides extra semantic meaning to the type. Types are typically denoted by $\sigma$, $\tau$. To say that term $e$ has type $\tau$ we use the following notation
\begin{align}
    e : \tau
\end{align}
Read this expression as "e has type tau". Function types are represented by the following notation
\begin{align}
    e: \sigma \to \tau
\end{align}
Functions with multiple arguments by chaining arrows. For example, here's the type of a function with 3 arguments.
\begin{align}
    e: \tau_1 \to \tau_2 \to \tau_3 \to \tau_4
\end{align}

The type of a list is denoted by
\begin{align}
    e: [\tau]
\end{align}


When evaluating types for a program, we make use of what is called a type environment, typically denoted by $\Gamma$. Under the hood, a type environment is simply a list of pairs $e: \tau$
\begin{align}
    \Gamma := [e_1: \tau_1,\dotsc,e_n: \tau_n]
\end{align}
Judgments are the process by a type environment is used to derive a new term-type assignment. A judgment is written as
\begin{align}
    \Gamma \vdash e: \tau
\end{align}
Read the above statement as "type environment Gamma implies that term e is of type tau".

\section{Typing Rules}

Typing rules are rules for logical inference that describe how a given type system assigns a type to every term. A program is said to be well typed if every term has an associated type. An inference rule is given by the following general form.

\begin{align}
    \inference[]
    {
     \textrm{premises}
    }
    {
     \textrm{conclusion}
    }
\end{align}

If all the premises above the line are fulfilled, the conclusions below the line follows by the typing rule. Let's look at a more concrete example.

\begin{align}
    \inference[]
    {
    \Gamma \vdash e_1 : \mathbb{R}, \quad \Gamma \vdash e_2 : \mathbb{R}
    }
    {
    \Gamma \vdash e_1 + e_2 : \mathbb{R}
    }
\end{align}

The trick to reading typing rules is to convert them into English. The typing rule above corresponds to the following English statement: 
\\ \\
"If expression $e_1$ has type $\mathbb{R}$ in type environment $\Gamma$, and if expression $e_2$ has type $\mathbb{R}$ in type environment $\Gamma$, then expression $e_1 + e_2$ has type $\mathbb{R}$ in type environment $\Gamma$." 
\\ \\
Using common sense, this typing rule is encoding the statement that the sum of two real numbers is a real number. A type system for a programming language is specified by a series of typing rules like the one displayed above. We shall see multiple examples of type systems in the following chapters.

\section{Exercises}

\begin{enumerate}
    \item We gave an example above of a typing rule for addition of real numbers. Construct typing rules for subtraction, multiplication, and division of real numbers.
\end{enumerate}
\end{document}